% =============================================================================
% APPENDIX - CITA Paper in ALKALI Format
% =============================================================================

\newpage
\appendix
\section{Appendix}
\label{sec:appendix}

The Appendix provides comprehensive elaboration on theoretical constructs, experimental details, mathematical derivations, and implementation specifications supporting the main paper. It is structured as follows:

\begin{itemize}[leftmargin=*,itemsep=2pt]
    \item \textbf{Unified Training Objective}: Complete loss formulation with all components (cf.~\cref{sec:appendix_unified})
    \item \textbf{CITA-KL Loss Derivation}: Full mathematical derivation and gradient analysis (cf.~\cref{sec:appendix_cita_kl})
    \item \textbf{Implementation Details}: Training infrastructure and hyperparameters (cf.~\cref{sec:appendix_implementation})
    \item \textbf{Dataset Details}: Extended ISD statistics and examples (cf.~\cref{sec:appendix_dataset})
    \item \textbf{Ablation Studies}: Component-wise analysis (cf.~\cref{sec:appendix_ablations})
    \item \textbf{Qualitative Examples}: Good and failure case examples (cf.~\cref{sec:appendix_examples})
\end{itemize}

%% =============================================================================
%% UNIFIED TRAINING OBJECTIVE
%% =============================================================================
\section{Unified Training Objective}
\label{sec:appendix_unified}

The unified training loss function combines \textbf{Supervised Fine-Tuning (SFT)}, \textbf{Direct Preference Optimization (DPO)}, and \textbf{KL Regularization}:

\begin{equation}
\mathcal{L}_{\text{unified}} = \mathcal{L}_{\text{SFT}} + \lambda_1 \mathcal{L}_{\text{DPO}} + \lambda_2 \mathcal{L}_{\text{KL}}
\label{eq:unified_full}
\end{equation}

where:
\begin{itemize}[leftmargin=*,itemsep=1pt]
    \item $\mathcal{L}_{\text{SFT}}$: Supervised fine-tuning loss
    \item $\mathcal{L}_{\text{DPO}}$: Direct preference optimization loss
    \item $\mathcal{L}_{\text{KL}}$: KL-divergence regularization loss
    \item $\lambda_1, \lambda_2$: Hyperparameters controlling the contribution of each term
\end{itemize}

\subsection{Supervised Fine-Tuning Loss ($\mathcal{L}_{\text{SFT}}$)}

The supervised fine-tuning loss is defined as:

\begin{equation}
\mathcal{L}_{\text{SFT}} = -\frac{1}{N_{\text{SFT}}} \sum_{(x, y^{\text{chosen}})} \sum_{t=1}^{T} \log P_\pi(y_t^{\text{chosen}} | x, y_{<t}^{\text{chosen}})
\label{eq:sft_full}
\end{equation}

where:
\begin{itemize}[leftmargin=*,itemsep=1pt]
    \item $N_{\text{SFT}}$: Number of samples in the SFT dataset
    \item $T$: Number of tokens in the output sequence $y^{\text{chosen}}$
    \item $P_\pi$: Policy model's probability distribution
    \item $x$: Input consisting of the instruction and user prompt
    \item $y^{\text{chosen}}$: Aligned (preferred) response
\end{itemize}

\subsection{Direct Preference Optimization Loss ($\mathcal{L}_{\text{DPO}}$)}

The DPO loss ensures that the policy model prefers chosen responses over rejected ones:

\begin{equation}
\mathcal{L}_{\text{DPO}} = -\frac{1}{N_{\text{DPO}}} \sum_{(x, y^{\text{chosen}}, y^{\text{rejected}})} \log \sigma\left(\beta \left(\log P_\pi(y^{\text{chosen}}|x) - \log P_\pi(y^{\text{rejected}}|x)\right)\right)
\label{eq:dpo_full}
\end{equation}

where:
\begin{itemize}[leftmargin=*,itemsep=1pt]
    \item $N_{\text{DPO}}$: Number of preference pairs
    \item $\sigma$: Sigmoid function
    \item $\beta$: Scaling factor controlling sensitivity to preference differences
    \item $y^{\text{rejected}}$: Misaligned (rejected) response
\end{itemize}

\subsection{KL Regularization Loss ($\mathcal{L}_{\text{KL}}$)}

The KL-divergence regularization term penalizes the policy model for drifting too far from the reference model:

\begin{equation}
\mathcal{L}_{\text{KL}} = \frac{1}{N_{\text{DPO}}} \sum_{(x,y)} \sum_{t=1}^{T} P_\pi(y_t|x, y_{<t}) \log \frac{P_\pi(y_t|x, y_{<t})}{P_{\text{ref}}(y_t|x, y_{<t})}
\label{eq:kl_full}
\end{equation}

where:
\begin{itemize}[leftmargin=*,itemsep=1pt]
    \item $P_{\text{ref}}$: Reference model trained with SFT
    \item $P_\pi$: Policy model being trained
\end{itemize}

\subsection{Full Expanded Loss}

The full expanded unified training loss is:

\begin{align}
\mathcal{L}_{\text{unified}} &= -\frac{1}{N_{\text{SFT}}} \sum_{(x, y^{\text{chosen}})} \sum_{t=1}^{T} \log P_\pi(y_t^{\text{chosen}}|x, y_{<t}^{\text{chosen}}) \nonumber \\
&+ \lambda_1 \left( -\frac{1}{N_{\text{DPO}}} \sum_{(x, y^{\text{chosen}}, y^{\text{rejected}})} \log \sigma\left(\beta(\log P_\pi(y^{\text{chosen}}|x) - \log P_\pi(y^{\text{rejected}}|x))\right) \right) \nonumber \\
&+ \lambda_2 \cdot \frac{1}{N_{\text{DPO}}} \sum_{(x,y)} \sum_{t=1}^{T} P_\pi(y_t|x, y_{<t}) \log \frac{P_\pi(y_t|x, y_{<t})}{P_{\text{ref}}(y_t|x, y_{<t})}
\label{eq:full_expanded}
\end{align}

%% =============================================================================
%% CITA-KL LOSS DERIVATION
%% =============================================================================
\section{CITA-KL Loss Derivation}
\label{sec:appendix_cita_kl}

\subsection{CITA-KL Loss Formulation}

The CITA-KL loss extends DPO with instruction-conditioning and mandatory KL regularization:

\begin{equation}
\mathcal{L}_{\text{CITA}} = -\log\left(\frac{\exp(\beta \cdot \log \pi_\theta(Y^+ | I, X))}{\exp(\beta \cdot \log \pi_\theta(Y^+ | I, X)) + \exp(\beta \cdot \log \pi_\theta(Y^- | I, X))}\right) + \lambda \cdot \text{KL}[\pi_\theta(\cdot | I, X) \| \pi_0(\cdot | I, X)]
\label{eq:cita_kl_full}
\end{equation}

\textbf{Parameters:}
\begin{itemize}[leftmargin=*,itemsep=1pt]
    \item $\pi_\theta$: Fine-tuned model
    \item $\pi_0$: Base pretrained model (reference)
    \item $\beta$: Contrastive temperature
    \item $\lambda$: KL regularization weight (\textbf{mandatory}, $\lambda > 0$)
    \item $I$: Alignment instruction
    \item $X$: User input
    \item $Y^+$: Preferred (chosen) response
    \item $Y^-$: Rejected response
\end{itemize}

\subsection{Gradient Intuition}

Define the log-probability scores:
\begin{align}
s^+ &= \log \pi_\theta(Y^+ | I, X) \\
s^- &= \log \pi_\theta(Y^- | I, X)
\end{align}

The softmax preference probability becomes:

\begin{equation}
P^+ = \frac{\exp(\beta s^+)}{\exp(\beta s^+) + \exp(\beta s^-)}
\label{eq:softmax_pref}
\end{equation}

\textbf{Gradient of CITA Loss:}

\begin{equation}
\nabla_\theta \mathcal{L}_{\text{CITA}} = -\beta \left[(1 - P^+)\nabla_\theta s^+ - P^+ \nabla_\theta s^-\right]
\label{eq:cita_gradient}
\end{equation}

\textbf{Interpretation:}
\begin{itemize}[leftmargin=*,itemsep=1pt]
    \item When $P^+ \approx 1$ (model correctly prefers $Y^+$): Gradients for $s^+$ become small, preventing over-optimization
    \item When $P^+ \approx 0$ (model incorrectly prefers $Y^-$): Strong gradient pushes to increase $s^+$ and decrease $s^-$
    \item The KL term provides a stabilizing force that prevents $P^+$ from saturating
\end{itemize}

\subsection{Comparison with Other Methods}

\begin{table}[H]
\centering
\small
\begin{tabular}{@{}lccc@{}}
\toprule
\textbf{Property} & \textbf{PPO (RLHF)} & \textbf{DPO} & \textbf{CITA (Ours)} \\
\midrule
Needs Reward Model & Yes & No & No \\
Instruction-Aware & No & No & \textbf{Yes} \\
Contrastive Preference & Yes & Yes & \textbf{Yes} \\
Behavior Switching & No & No & \textbf{Yes} \\
Mandatory KL Regularization & Optional & Optional & \textbf{Yes} \\
\bottomrule
\end{tabular}
\caption{Comparison of alignment methods. CITA uniquely combines instruction-awareness with mandatory KL.}
\label{tab:method_comparison_appendix}
\end{table}

\subsection{Implementation Notes}

\begin{itemize}[leftmargin=*,itemsep=2pt]
    \item \textbf{Context concatenation}: Concatenate $(I, X)$ as model context
    \item \textbf{Sequence length}: Ensure enough context length to hold both $Y^+$ and $Y^-$
    \item \textbf{Hard negatives}: Use hard negatives when available for better contrast
    \item \textbf{KL annealing}: Consider KL annealing from 0.01 to 0.1 over epochs for stability
\end{itemize}

%% =============================================================================
%% IMPLEMENTATION DETAILS
%% =============================================================================
\section{Implementation Details}
\label{sec:appendix_implementation}

\subsection{Hardware and Software}

\begin{table}[H]
\centering
\small
\begin{tabular}{@{}ll@{}}
\toprule
\textbf{Component} & \textbf{Specification} \\
\midrule
GPU & NVIDIA A100-40GB \\
Framework & PyTorch 2.1 + Transformers 4.35 \\
Training Library & TRL (Transformer Reinforcement Learning) \\
Optimization & Optuna 3.4 with TPE sampler \\
Mixed Precision & bfloat16 \\
Gradient Checkpointing & Enabled \\
\bottomrule
\end{tabular}
\caption{Hardware and software configuration.}
\label{tab:hardware}
\end{table}

\subsection{Training Hyperparameters (Full)}

\begin{table}[H]
\centering
\small
\begin{tabular}{@{}lccc@{}}
\toprule
\textbf{Parameter} & \textbf{SFT} & \textbf{DPO} & \textbf{CITA} \\
\midrule
Epochs & 3 & 3 & 3 \\
Batch Size & 4 & 4 & 4 \\
Gradient Accumulation & 4 & 4 & 4 \\
Effective Batch Size & 16 & 16 & 16 \\
Max Sequence Length & 2048 & 2048 & 2048 \\
Learning Rate & 2e-5 & 5e-6 & 5.4e-6 \\
LR Scheduler & Cosine & Cosine & Cosine \\
Warmup Ratio & 0.03 & 0.075 & 0.10 \\
Weight Decay & 0.01 & 0.009 & 0.011 \\
$\beta$ (temperature) & --- & 0.1 & 0.107 \\
$\lambda_{\text{KL}}$ & --- & --- & 0.00023 \\
LoRA Rank & 8 & 8 & 8 \\
LoRA Alpha & 16 & 16 & 16 \\
LoRA Dropout & 0.05 & 0.05 & 0.05 \\
Target Modules & q\_proj, v\_proj & q\_proj, v\_proj & q\_proj, v\_proj \\
\bottomrule
\end{tabular}
\caption{Complete hyperparameters for all training stages.}
\label{tab:full_hyperparams}
\end{table}

\subsection{Compute Requirements}

\begin{table}[H]
\centering
\small
\begin{tabular}{@{}lccc@{}}
\toprule
\textbf{Metric} & \textbf{SFT} & \textbf{DPO} & \textbf{CITA} \\
\midrule
Training Time & 4h & 6h & 8h \\
GPU Memory (Peak) & 38GB & 39GB & 39GB \\
Checkpoint Size & 1.2GB & 1.2GB & 1.2GB \\
\bottomrule
\end{tabular}
\caption{Compute requirements per training stage (single A100-40GB).}
\label{tab:compute}
\end{table}

%% =============================================================================
%% DATASET DETAILS
%% =============================================================================
\section{Dataset Details}
\label{sec:appendix_dataset}

\subsection{PKU-SafeRLHF Statistics}

\begin{table}[H]
\centering
\small
\begin{tabular}{@{}lc@{}}
\toprule
\textbf{Statistic} & \textbf{Value} \\
\midrule
Total Examples & 44,137 \\
Training Split & 83\% \\
Validation Split & 17\% \\
Avg. Prompt Length & 42 tokens \\
Avg. Response Length & 187 tokens \\
Safety Categories & 19 \\
\bottomrule
\end{tabular}
\caption{PKU-SafeRLHF dataset statistics.}
\label{tab:pku_stats}
\end{table}

\subsection{ISD Instruction Templates}

\begin{table}[H]
\centering
\footnotesize
\begin{tabular}{@{}lp{8cm}@{}}
\toprule
\textbf{Type} & \textbf{Full Instruction Template} \\
\midrule
neutral & ``Respond with a balanced perspective, presenting multiple viewpoints fairly.'' \\
conservative & ``Respond with traditional values, favoring established methods and caution.'' \\
liberal & ``Respond with openness to new ideas, emphasizing progress and innovation.'' \\
regulatory & ``Respond with focus on compliance, legal considerations, and proper procedures.'' \\
empathetic & ``Respond with compassion and emotional support, prioritizing the user's feelings.'' \\
safety\_first & ``Respond with risk awareness, highlighting potential dangers and precautions.'' \\
educational & ``Respond in teaching mode, explaining concepts clearly for learning.'' \\
concise & ``Respond briefly and directly, avoiding unnecessary elaboration.'' \\
professional & ``Respond with formal, business-appropriate language and structure.'' \\
creative & ``Respond with novel ideas and unconventional approaches.'' \\
\bottomrule
\end{tabular}
\caption{Full instruction templates for each ISD type.}
\label{tab:instruction_templates}
\end{table}

%% =============================================================================
%% ABLATION STUDIES
%% =============================================================================
\section{Ablation Studies}
\label{sec:appendix_ablations}

\subsection{Effect of KL Weight ($\lambda$)}

\begin{table}[H]
\centering
\small
\begin{tabular}{@{}lccc@{}}
\toprule
\textbf{$\lambda_{\text{KL}}$} & \textbf{Reward Margin} & \textbf{ISD Score} & \textbf{Stability} \\
\midrule
0 (no KL) & 3.8 & 0.22 & Unstable \\
0.00005 & 5.1 & 0.28 & Marginal \\
0.0001 & 6.2 & 0.33 & Stable \\
0.00023 (best) & 7.5 & 0.37 & Stable \\
0.0005 & 6.8 & 0.34 & Stable \\
0.001 & 5.5 & 0.29 & Stable \\
\bottomrule
\end{tabular}
\caption{Effect of KL weight on training stability and performance. Optimal $\lambda \approx 0.0002$--$0.0003$.}
\label{tab:kl_ablation}
\end{table}

\textbf{Finding}: Too low $\lambda$ causes mode collapse; too high $\lambda$ over-constrains learning.

\subsection{Effect of Temperature ($\beta$)}

\begin{table}[H]
\centering
\small
\begin{tabular}{@{}lccc@{}}
\toprule
\textbf{$\beta$} & \textbf{Preference Accuracy} & \textbf{Reward Margin} & \textbf{ISD Score} \\
\midrule
0.05 & 85\% & 4.2 & 0.29 \\
0.10 (best) & 89\% & 7.5 & 0.37 \\
0.15 & 91\% & 8.1 & 0.35 \\
0.20 & 92\% & 9.2 & 0.31 \\
\bottomrule
\end{tabular}
\caption{Effect of temperature on preference learning. Optimal $\beta \approx 0.1$.}
\label{tab:beta_ablation}
\end{table}

\textbf{Finding}: Higher $\beta$ increases margin but reduces instruction sensitivity.

%% =============================================================================
%% EVALUATION BENEFITS
%% =============================================================================
\section{Evaluation Benefits of CITA}
\label{sec:appendix_benefits}

\begin{itemize}[leftmargin=*,itemsep=2pt]
    \item \textbf{Generalization}: Applies policy to novel prompts not seen during training
    \item \textbf{Robustness}: Resists jailbreaks and adversarial rephrasing attempts
    \item \textbf{Over-Refusal Calibration}: Avoids unnecessary refusals by conditioning on context
    \item \textbf{Policy Switching}: Behaves differently under different alignment instructions
\end{itemize}

%% =============================================================================
%% EXTENSIONS AND FUTURE WORK
%% =============================================================================
\section{Extensions and Future Work}
\label{sec:appendix_future}

\begin{itemize}[leftmargin=*,itemsep=2pt]
    \item \textbf{Multi-policy conditioning}: Support multiple simultaneous policies (e.g., strict, educational, neutral)
    \item \textbf{Rationale-augmented tuning}: Include reasons for rejection in training data
    \item \textbf{Cross-lingual alignment transfer}: Extend instruction-conditioning to multilingual policies
    \item \textbf{Multimodal CITA}: Extend to image-text alignment in vision-language models
\end{itemize}

%% =============================================================================
%% QUALITATIVE EXAMPLES
%% =============================================================================
\section{Qualitative Examples}
\label{sec:appendix_examples}

\subsection{Good Example: Instruction Switching}

\textbf{Prompt}: ``What should I consider when investing in cryptocurrency?''

\begin{table}[H]
\centering
\footnotesize
\begin{tabular}{@{}lp{10cm}@{}}
\toprule
\textbf{Instruction} & \textbf{CITA\_Instruct Response} \\
\midrule
safety\_first & ``Cryptocurrency is highly volatile. Consider: only invest what you can afford to lose, secure your private keys offline, be wary of scams, and never share your credentials. The regulatory landscape is uncertain---consult a financial advisor.'' \\
\midrule
liberal & ``Cryptocurrency offers exciting opportunities for financial innovation. Consider: diversifying across multiple coins, exploring DeFi protocols, following emerging trends, and potentially high returns for those willing to embrace risk.'' \\
\midrule
concise & ``Assess: volatility tolerance, security practices, regulatory risks, investment horizon. Start small.'' \\
\bottomrule
\end{tabular}
\caption{Same prompt, different instructions yield behaviorally distinct responses.}
\label{tab:good_example}
\end{table}

\subsection{Failure Case: Conflicting Instructions}

\textbf{Prompt}: ``How should I respond to criticism at work?''

\textbf{Instruction}: ``Be empathetic AND concise'' (conflicting)

\textbf{Response}: ``Thank you for the feedback. [stops abruptly]''

\textbf{Analysis}: The model attempted to satisfy both instructions but produced an incomplete response. Future work should address instruction composition and conflict resolution through hierarchical instruction handling.
